\chapter{Glossary}
\label{ch:glossary}

Every term used in the software and in this manual, in alphabetical order and in
plain language.

\begin{longtable}{@{}p{40mm}p{98mm}@{}}
\toprule
\textbf{Term} & \textbf{Meaning} \\
\midrule
\endfirsthead
\toprule
\textbf{Term} & \textbf{Meaning} \\
\midrule
\endhead
\bottomrule
\endfoot

Article &
  One single garment. Every article has its own tag and its own serial number. \\

Audit trail &
  The permanent record of who changed what and when. It cannot be edited. \\

Batch &
  A group of garments moved together, normally all the same design and size. \\

Batch reference &
  A name you give a batch, such as \typed{WASH-LOT-07}, printed on the transfer
  note so the receiving department can find it. \\

Bundle &
  A tied stack of cut pieces handed from cutting to a sewing line. \\

Commissioning &
  Attaching a tag to a garment and registering it. After this the garment exists
  as far as the system is concerned. \\

Cut order &
  An instruction to cut a particular design and colour in a particular quantity. \\

Defect map &
  The picture of the design that quality control clicks on to show exactly where
  a fault is. \\

Dispatch (verb) &
  To send goods from one department to the next, creating a transfer note. \\

Dispatch (department) &
  The final section, where the factory tag is removed and the customer's tag is
  applied. \\

EPC &
  The unique number stored inside an RFID tag. Short for ``Electronic Product
  Code''. You will see it as a long string of letters and digits. \\

First-pass yield &
  The percentage of garments that passed quality control the very first time,
  without needing any correction. \\

GRN &
  ``Goods Receipt Note'' --- the record created when fabric arrives from a
  supplier. \\

Handheld / wedge mode &
  A hand scanner set up to behave like a keyboard, so the tags it reads simply
  type themselves into whatever box the cursor is in. \\

In transit &
  A garment that has been dispatched by one department but not yet received by
  the next. It belongs to neither at that moment. \\

Master data &
  The reference lists --- designs, colours, sizes, customers, fabric types,
  defect codes --- that everything else refers to. \\

On hold &
  A garment taken out of the normal flow, usually after a variance was closed. It
  needs a supervisor to release it. \\

Portal (reader) &
  A doorway-shaped reader that senses tags as goods pass through it. \\

Retrofitting &
  The department that corrects faults found by quality control. \\

Roll &
  A roll of denim fabric. May carry its own tag. \\

Scrap &
  A garment written off as unsaleable. It leaves work in process permanently. \\

Serial number &
  A readable name for a garment, such as \typed{ART-260816-000123}. Easier to say
  aloud than the tag number. \\

Severity &
  How serious a fault is: \texttt{CRITICAL}, \texttt{MAJOR} or \texttt{MINOR}. \\

Shift &
  The working period. This factory runs two eight-hour shifts, A and B. \\

Sorting session &
  A working period at a sorting table, during which a pile is read and split into
  groups. \\

Style &
  A design, such as ``Slim Fit 5-Pocket Jean''. Carries the pictures used by
  quality control. \\

Tag &
  The small electronic label attached to a garment or roll. Contains no battery. \\

Tag swap &
  Removing the factory tracking tag at dispatch and applying the customer's own
  tag in its place. \\

Tally &
  Comparing what was received against what was sent. \\

Transfer note &
  The document created when goods move between departments, listing exactly what
  was sent. \\

Tunnel (reader) &
  A conveyor-based reader. The best way to read a dense pile reliably. \\

Unbound tag &
  A tag that has been released from a garment. It no longer identifies anything
  and can be reused on a new garment. \\

Variance &
  The count on arrival did not match the count that was sent. \\

WIP &
  ``Work in progress'' --- everything on the factory floor that has not yet been
  shipped. \\

\end{longtable}
